% !TeX program = pdflatex
\documentclass[11pt,a4paper]{amsart}

\newcommand{\notelanguage}{finnish}
% Shared preamble for course notes and exercise sets.

\usepackage[T1]{fontenc}
\usepackage{lmodern}
\usepackage[english]{babel}

\usepackage[a4paper,margin=30mm]{geometry}
\usepackage{microtype}
\usepackage{graphicx}
\usepackage{enumitem}

\usepackage{amsmath}
\usepackage{amssymb}
\usepackage{amsthm}
\usepackage{mathtools}
\usepackage{aliascnt}

\usepackage[hidelinks,pdfusetitle]{hyperref}

\numberwithin{equation}{section}
\setlist{itemsep=0.25em,topsep=0.5em}

\theoremstyle{plain}
\newtheorem{theorem}{Theorem}[section]
\newaliascnt{lemma}{theorem}
\newtheorem{lemma}[lemma]{Lemma}
\aliascntresetthe{lemma}
\newaliascnt{proposition}{theorem}
\newtheorem{proposition}[proposition]{Proposition}
\aliascntresetthe{proposition}
\newaliascnt{corollary}{theorem}
\newtheorem{corollary}[corollary]{Corollary}
\aliascntresetthe{corollary}

\theoremstyle{definition}
\newaliascnt{definition}{theorem}
\newtheorem{definition}[definition]{Definition}
\aliascntresetthe{definition}
\newaliascnt{example}{theorem}
\newtheorem{example}[example]{Example}
\aliascntresetthe{example}
\newaliascnt{problem}{theorem}
\newtheorem{problem}[problem]{Problem}
\aliascntresetthe{problem}

\theoremstyle{remark}
\newaliascnt{remark}{theorem}
\newtheorem{remark}[remark]{Remark}
\aliascntresetthe{remark}

\usepackage[nameinlink,noabbrev]{cleveref}

\crefname{theorem}{theorem}{theorems}
\crefname{lemma}{lemma}{lemmas}
\crefname{proposition}{proposition}{propositions}
\crefname{corollary}{corollary}{corollaries}
\crefname{definition}{definition}{definitions}
\crefname{example}{example}{examples}
\crefname{problem}{problem}{problems}
\crefname{remark}{remark}{remarks}

\DeclareMathOperator{\im}{im}
\DeclareMathOperator{\rank}{rank}
\DeclarePairedDelimiter{\abs}{\lvert}{\rvert}
\DeclarePairedDelimiter{\norm}{\lVert}{\rVert}
\DeclarePairedDelimiter{\set}{\lbrace}{\rbrace}


\usepackage{tikz}

\title{Analyysi I\\Ryhmätyö 4 ratkaisut}
\author{}
\date{}

\begin{document}

\exerciseheading{Analyysi I}{Ryhmätyö 4 ratkaisut}

\begin{exercise}
  Tutkitaan lukujono \((y_n)\), missä \(y_n = \frac{(n+1)^2}{n^2+1}\)
  kun \(n = 1, 2, 3, \ldots\). Haluamme selvittää lukujonosta seuraavat asiat:
  \begin{itemize}[nosep]
    \item Lukujonon raja-arvo.
    \item Onko lukujono vähenevä tai kasvava.
    \item \(\inf\{y_n : n = 1, 2, 3, \ldots\}\).
    \item \(\sup\{y_n : n = 1, 2, 3, \ldots\}\).
    \item Suppeneeko lukujono.
  \end{itemize}
  Missä järjestyksessä asiat kannattaa selvittää? Valitse sopiva järjestys ja
  tutki väitteet. Voiko väitteitä perustella eri tavoilla?
\end{exercise}

\begin{solution}
  Selvitetään ensin raja-arvo ja suppeneminen. Koska
  \[
    y_n = 1 + \frac{2n}{n^2+1},
    \qquad 0 < y_n-1 \leq \frac{2}{n} \to 0,
  \]
  lukujono suppenee ja sen raja-arvo on \(1\). Lisäksi
  \[
    y_{n+1}-y_n
    = \frac{-2(n^2+n-1)}{(n^2+1)((n+1)^2+1)} < 0
    \qquad(n\geq 1),
  \]
  joten lukujono on aidosti vähenevä. Näin ollen
  \(\sup\{y_n:n\geq1\}=y_1=2\). Koska kaikki termit ovat suurempia
  kuin \(1\) ja lähestyvät sitä, \(\inf\{y_n:n\geq1\}=1\).

  Vaihtoehtoisesti voidaan ensin osoittaa jono väheneväksi ja alhaalta
  rajoitetuksi luvulla \(1\). Tällöin suppeneminen seuraa monotonisen
  lukujonon suppenemislauseesta, ja raja-arvo saadaan jakamalla
  lausekkeen \(\frac{(n+1)^2}{n^2+1}\) osoittaja ja nimittäjä luvulla \(n^2\).
\end{solution}

\begin{exercise}
  \leavevmode\par
  \begin{enumerate}[label=(\arabic*), nosep]
    \item Olkoon \((x_n)\) lukujono, joka suppenee kohti lukua \(x\).
      Osoita, että jokainen lukujonon \((x_n)\) osajono suppenee kohti
      lukua \(x\).
    \item Olkoon \((x_{n_i})\) lukujonon \((x_n)\) osajono. Oletetaan,
      että \((x_{n_i})\) hajaantuu. Osoita, että \((x_n)\) hajaantuu.
    \item Olkoot \((x_{n_j})\) ja \((x_{n_k})\) lukujonon \((x_n)\)
      osajonoja. Oletetaan, että \(x_{n_j} \to a\), \(x_{n_k} \to b\)
      ja \(a \neq b\). Osoita, että lukujono \((x_n)\) hajaantuu.
  \end{enumerate}
\end{exercise}

\begin{solution}
  \leavevmode\par
  \begin{enumerate}[label=(\arabic*), nosep]
    \item Olkoon \((x_{n_i})\) osajono ja \(\epsilon>0\). Koska
      \(x_n\to x\), on olemassa \(N\in\NNone\) siten, että
      \(\abs{x_n-x}<\epsilon\) kaikilla \(n > N\).
      Osajonon indeksit ovat aidosti kasvavia, joten \(n_i\geq i\).
      Kun \(i > N\), siis myös \(n_i > N\), ja näin
      \(\abs{x_{n_i}-x}<\epsilon\). Siis \(x_{n_i}\to x\).
    \item Väite on kohdan (1) kontrapositio. Jos \((x_n)\) suppenisi,
      myös \((x_{n_i})\) suppenisi kohdan
      (1) nojalla. Tämä on ristiriidassa oletuksen kanssa, joten
      \((x_n)\) hajaantuu.
    \item Jos \(x_n\to x\), niin kohdan (1) nojalla molemmat osajonot
      suppenisivat kohti lukua \(x\). Raja-arvon yksikäsitteisyydestä
      seuraisi \(a=x=b\), mikä on ristiriita. Siis \((x_n)\) hajaantuu.
  \end{enumerate}
\end{solution}

\end{document}
